\chapter{Methodological Framework}

\section{The Epistemological Horizon of Comparative Philology}
Traditional comparative linguistics successfully reconstructs language families back to approximately 6,000 to 8,000 BP. However, beyond the 10,000-year mark, the methodology faces significant challenges. This study addresses the '10,000-year decay limit' not by ignoring it, but by shifting the analytical focus from high-volatility lexical items to structural invariants that exhibit higher resistance to temporal decay.

\section{Addressing the Constraints of Long-Range Comparison}
Mainstream historical linguistics has rightfully viewed macro-family reconstructions (e.g., Nostratic, Eurasiatic) with extreme skepticism. As established by Campbell and Poser (2008), the primary failure mode of such proposals is the 'look-alike' fallacy, where coincidental phonetic similarities are mistaken for genetic relationship. Furthermore, Nichols (1992) has demonstrated that structural typological features (such as SOV order and agglutination) are subject to horizontal diffusion across unrelated languages through areal contact.

\section{Structural Modeling and Signal Invariance}
To mitigate these risks, this study utilizes a multi-disciplinary framework that triangulates linguistic structuralism with archaeological stratigraphy and computational validation. We distinguish between 'genetic invariants' and 'typological accidents' by requiring:
\begin{enumerate}
    \item \textbf{Non-Linear Phonetic Triangulation:} Every reconstruction must follow a systematic sound law (The Glottalic Shift) that produces distinct reflexes in the three primary branches.
    \item \textbf{Zero-Shot Predictive Modeling:} The validity of the proposed syntax is tested on its ability to predict structural properties in unseen datasets, specifically the Indus Valley Script, compared against a structural null model.
\end{enumerate}

\section{Primary Data Anchoring}
The methodology utilizes primary attested data from archaic literary and archaeological strata, rather than second-order theoretical reconstructions. This study is anchored in:
\begin{enumerate}
    \item \textbf{Vedic Sanskrit (c. 1500 BCE):} The Western Indo-European anchor.
    \item \textbf{Old Tamil (c. 300 BCE):} The Southern Dravidian anchor.
    \item \textbf{Old Tibetan (c. 700 CE):} The Eastern Sino-Tibetan anchor.
\end{enumerate}

\section{Probability of Convergence}
The probability of accidental alignment across independent data streams is modeled through Bayesian convergence. We define the unified signal as the intersection of independent linguistic ($L$), computational ($C$), and archaeological ($A$) evidence, where the joint probability $P(L \cap C \cap A)$ falls significantly below the threshold for accidental coincidence.

